% ============================================================================
% Chapitre 9 — Résultats, discussion et perspectives
% ============================================================================
\chapter{Résultats, discussion et perspectives}

\section*{Introduction}

Ce chapitre confronte les résultats obtenus aux objectifs initiaux, analyse
les forces et les limites de la plateforme de manière objective, et décrit
les perspectives d'évolution identifiées au fil du développement.

\section{Résultats obtenus}

\subsection{Fonctionnalités livrées}

L'ensemble des objectifs définis au chapitre introductif ont été atteints.
Le tableau~\ref{tab:obj-results} dresse la matrice de conformité :

\begin{table}[H]
\centering
\caption{Matrice objectifs initiaux -- résultats obtenus}
\label{tab:obj-results}
\begin{tabularx}{\textwidth}{lXl}
\toprule
\textbf{Objectif} & \textbf{Résultat} & \textbf{Statut} \\
\midrule
O1 : Dashboard answers-first
  & Implémenté : statut global, problèmes, actions recommandées
  & \textcolor{SuccessGreen}{\textbf{OK}} \\
O2 : Corrélation des changements
  & Moteur complet : scoring temporel + LLM narrative
  & \textcolor{SuccessGreen}{\textbf{OK}} \\
O3 : Explication LLM local
  & Ollama + Llama 3.1 8B, prompts non-DevOps
  & \textcolor{SuccessGreen}{\textbf{OK}} \\
O4 : Prévision Holt-Winters 24h
  & Implémentée : 24 points, time-to-critical, narratif
  & \textcolor{SuccessGreen}{\textbf{OK}} \\
O5 : Architecture microservices
  & 19 services, NATS JetStream, 8 streams
  & \textcolor{SuccessGreen}{\textbf{OK}} \\
O6 : Sécurité bout en bout
  & mTLS, JWT, RLS, AES-256-GCM, audit log
  & \textcolor{SuccessGreen}{\textbf{OK}} \\
O7 : Déploiement $\leq$ 10 min
  & 6 étapes, environ 8 min hors téléchargement modèles
  & \textcolor{SuccessGreen}{\textbf{OK}} \\
\bottomrule
\end{tabularx}
\end{table}

\subsection{Métriques de la base de code}

\begin{table}[H]
\centering
\caption{Indicateurs quantitatifs de la base de code \caimp{}}
\label{tab:codebase-metrics}
\begin{tabularx}{\textwidth}{lX}
\toprule
\textbf{Indicateur} & \textbf{Valeur} \\
\midrule
Lignes de code total (hors tests, hors commentaires)
  & $\sim$ 52~000 \\
Services Python & 11 (admin-api, query-api, ai-worker, alert-engine,
  forecast-engine, correlation-engine, change-tracker, ws-gateway,
  ai-orchestrator, splunk-hec-bridge, agent) \\
Services Go     & 1 (telemetry-writer) \\
Composants React / TypeScript & 26 composants, 6 pages \\
Migrations SQL  & 14 migrations (extensions → isolations RLS) \\
Tables PostgreSQL & 22 tables \\
Streams NATS JetStream & 8 streams \\
Tests unitaires & 34 tests pytest \\
Couverture tests critiques & $\sim$74~\% (détecteurs + forecaster + scoring) \\
\bottomrule
\end{tabularx}
\end{table}

\section{Forces du projet}

\begin{description}[leftmargin=2cm]
  \item[Architecture découplée] L'utilisation de NATS JetStream comme bus
    d'événements garantit qu'aucun service ne dépend directement d'un autre.
    L'ajout d'un nouveau consommateur d'anomalies (par exemple, un service de
    remédiation automatique) ne nécessite aucune modification des services
    existants.
  \item[Explication en langage naturel ciblée] Le reformatage des prompts
    LLM pour un public non-DevOps est une différenciation forte. Les tests
    qualitatifs montrent que les explications générées sont directement
    compréhensibles par des développeurs sans formation en exploitation système.
  \item[Corrélation de changements] Le moteur de corrélation automatique
    est absent de toutes les solutions open source comparées dans l'état de l'art.
    En associant chaque anomalie aux changements récents, il réduit le temps
    de diagnostic de plusieurs dizaines de minutes à quelques secondes.
  \item[Zéro dépendance cloud] L'ensemble du traitement IA s'effectue localement.
    Les données ne quittent jamais le serveur de l'opérateur, ce qui satisfait
    les exigences de confidentialité les plus strictes.
  \item[Portabilité] Le déploiement Docker Compose fonctionne identiquement
    sur Linux, macOS et Windows avec Docker Desktop.
\end{description}

\section{Limites identifiées}

Une analyse honnête du projet révèle plusieurs limitations à documenter :

\begin{description}[leftmargin=2cm]
  \item[Absence de fenêtre saisonnière dans Holt-Winters] Le modèle double
    exponential ne capture pas les variations cycliques (pic CPU quotidien,
    trafic hebdomadaire). L'extension à Holt-Winters triple (avec composante
    saisonnière) améliorerait la précision des prévisions à 24h.
  \item[Qualité des explications LLM dépendante du matériel] Sur une machine
    sans GPU, l'inférence Llama 3.1 8B peut prendre 15 à 30 secondes.
    Ce délai est acceptable pour une explication asynchrone, mais perceptible
    pour l'utilisateur.
  \item[Scalabilité du moteur de corrélation] La requête SQL sur \texttt{change\_events}
    dans une fenêtre de 30 minutes est efficace jusqu'à quelques milliers
    de changements simultanés. Au-delà (infrastructure de très grande taille),
    une approche par index temporel dédié s'imposerait.
  \item[Support Windows pour l'agent] L'agent Python utilise \texttt{psutil}
    qui supporte Windows, mais les reporters de changements (auth.log, dpkg.log)
    sont spécifiques Linux.
  \item[Absence de remédiation automatisée] \caimp{} produit des recommandations
    d'action mais n'exécute aucune correction automatique. Cette limitation
    est volontaire (principe de précaution) mais réductrice pour les équipes
    cherchant une plateforme autonome.
\end{description}

\section{Perspectives d'évolution}

\subsection{Court terme (3--6 mois)}

\begin{itemize}
  \item \textbf{Holt-Winters triple} : ajout de la composante saisonnière
    pour capturer les variations journalières et hebdomadaires ;
  \item \textbf{Support Windows} pour les reporters de changements de l'agent ;
  \item \textbf{Alertes Slack / Teams natifs} : intégration directe sans
    webhook générique ;
  \item \textbf{Fine-tuning du modèle} sur les incidents de l'organisation :
    utilisation du \textit{feedback} thumbs-up/down pour constituer un
    dataset d'entraînement.
\end{itemize}

\subsection{Moyen terme (6--18 mois)}

\begin{itemize}
  \item \textbf{Corrélation de traces} : intégration d'OpenTelemetry traces
    dans le moteur de corrélation pour associer les anomalies aux traces
    de requêtes défectueuses ;
  \item \textbf{Remédiation automatisée assistée} : exécution de scripts de
    remédiation pré-approuvés (runbooks exécutables) avec confirmation de
    l'opérateur ;
  \item \textbf{Mode SaaS hébergé} : offre commerciale hébergée pour les
    équipes ne souhaitant pas administrer leur propre déploiement, avec
    isolation Kubernetes par client ;
  \item \textbf{Tableau de bord mobile} : application React Native avec
    notifications push pour les alertes critiques.
\end{itemize}

\subsection{Vision long terme}

La vision à long terme de \caimp{} est de devenir l'ingénieur SRE
(\textit{Site Reliability Engineer}) virtuel des petites équipes --- un
système capable non seulement d'expliquer ce qui s'est passé, mais
d'anticiper les problèmes futurs et de les corriger de manière autonome,
tout en maintenant l'opérateur informé et en contrôle.

\section*{Conclusion}

Les résultats obtenus démontrent que les objectifs fixés ont été pleinement
atteints dans le délai d'un mois imparti. La plateforme \caimp{} livre une
proposition de valeur unique dans le paysage de la supervision open source :
des explications en langage naturel, une corrélation de changements automatique
et un déploiement accessible à tous. Les limitations identifiées constituent
un programme clair pour les évolutions futures.
